\unit{Adding angular momenta, Clebsch--Gordan and 3-j symbols}
\unitref{Week 8; 2025 Q1; Atomic Assignment 4 Q1}
\subsection{What is being added?}
An electron has orbital angular momentum
\(\mathbf L\) and spin \(\mathbf S\);
an atom may have nuclear angular momentum
\(\mathbf I\). The sum of two angular
momentum vectors is again an angular
momentum operator:
\(\mathbf J=\mathbf J_1+\mathbf J_2\).
But a simultaneous eigenstate of
\(J_1^2,J_2^2,J^2,J_z\) is not generally
a product of separate \(J_{1z}\) and
\(J_{2z}\) eigenstates.

Angular momentum eigenvalues are
\[
J^2\ket{jm}=\hbar^2j(j+1)\ket{jm},
\qquad J_z\ket{jm}=\hbar m\ket{jm}.
\]
For \(j_1,j_2\), the allowed total \(j\) values are
\[
\boxed{j=|j_1-j_2|,\ |j_1-j_2|+1,\ldots,j_1+j_2.}
\]
The total projection is \(m=m_1+m_2\).
That is the \emph{angular-momentum addition rule}.

\subsection{What does a Clebsch--Gordan coefficient say?}
In the product basis,
\(\ket{j_1m_1}\ket{j_2m_2}\) describes
known projections of the two components.
In the coupled basis \(\ket{JM}\),
the magnitude of the \emph{total}
angular momentum is known.
The relation is a change of basis:
\[
\boxed{\ket{JM}=
\sum_{m_1m_2}
\braket{j_1m_1,j_2m_2}{JM}
\ket{j_1m_1}\ket{j_2m_2}.}
\]
The number in angle brackets is a
Clebsch--Gordan (CG) coefficient.
It is a probability \emph{amplitude}
for a particular uncoupled pair, not
a physical coupling energy or a probability by itself.

\begin{worked}{Two spin-\(\tfrac12\) particles from scratch}
Each particle has \(\ket\uparrow\) or
\(\ket\downarrow\). Addition allows
\(J=1\) (triplet) and \(J=0\) (singlet).
The coupled states are
\[
\begin{aligned}
\ket{1,1}&=\ket{\uparrow\uparrow},\\
\ket{1,0}&=\frac{\ket{\uparrow\downarrow}
+\ket{\downarrow\uparrow}}{\sqrt2},\\
\ket{1,-1}&=\ket{\downarrow\downarrow},\\
\ket{0,0}&=\frac{\ket{\uparrow\downarrow}
-\ket{\downarrow\uparrow}}{\sqrt2}.
\end{aligned}
\]
Thus
\(\langle\tfrac12,\tfrac12;
\tfrac12,-\tfrac12|1,0\rangle=1/\sqrt2\).
The same product basis component
has overlap \(1/\sqrt2\) with the singlet.
Triplet and singlet states are
orthogonal, and the four coupled
states span the same four-dimensional
Hilbert space as the four product states.
\end{worked}

\subsection{Why introduce a 3-j symbol?}
A Wigner 3-j symbol is a conventional
normalised version of a CG coefficient:
\[
\boxed{
\begin{pmatrix}j_1&j_2&J\\m_1&m_2&-M\end{pmatrix}
=\frac{(-1)^{j_1-j_2+M}}
{\sqrt{2J+1}}
\langle j_1m_1,j_2m_2|JM\rangle.}
\]
It is useful because its symmetry
properties make selection rules and
angular integrations compact.
Before calculating a value, test:
\[
|j_1-j_2|\le j_3\le j_1+j_2,\qquad
m_1+m_2+m_3=0,\qquad |m_i|\le j_i.
\]
If a condition fails, the symbol is zero.
This often answers a whole transition
question before numerical evaluation.

\begin{worked}{A concrete \(3j\) calculation}
Adding zero angular momentum to \(j=1\)
does nothing:
\[
\langle 0,0;1,0|1,0\rangle=1.
\]
Use the relation above with
\(j_1=0,j_2=1,J=1,M=0\):
\[
\begin{pmatrix}0&1&1\\0&0&0\end{pmatrix}
=\frac{(-1)^{0-1+0}}{\sqrt3}(1)
=\boxed{-\frac1{\sqrt3}}.
\]
There is no need to memorise the
full factorial formula for this special case.
For less obvious values, consult the
CG/3-j tables provided in the course
or use a symbolic-algebra program.
\end{worked}

\subsection{Wigner--Eckart: separate geometry from physics}
A spherical tensor operator \(T_q^{(k)}\)
has rank \(k\) and component \(q=-k,\ldots,k\).
Electric dipole is rank 1; electric
quadrupole is rank 2.
For a chosen reduced-element convention,
\[
\boxed{
\mel{j'm'}{T_q^{(k)}}{jm}
=(-1)^{j'-m'}
\begin{pmatrix}j'&k&j\\-m'&q&m\end{pmatrix}
\langle j'\Vert T^{(k)}\Vert j\rangle.}
\]
The 3-j symbol gives the allowed
\emph{geometry}, including the projection
condition \(m'=m+q\).
The reduced matrix element contains
the specific radial/atomic physics but
is independent of \(m,m',q\).
Two transitions within the same
multiplet can therefore have
different strengths entirely because
their 3-j coefficients differ.

\subsection{A transition example with \(J=0\to J'=1\)}
An E1 operator has \(k=1\). Begin
in \(J=0,m=0\). The triangle rule permits
\(J'=1\). If the light has tensor
component \(q=+1\), then \(m'=1\)
because \(m'=m+q\).
For \(q=0\), the only final projection is
\(m'=0\). In both cases the absolute
angular coefficient is \(1/\sqrt3\)
in the displayed Wigner--Eckart convention.
Whether the actual transition occurs
still depends on the \emph{reduced}
dipole element and parity.

\subsection{Why parity is an independent test}
Under spatial inversion,
\(\mathbf r\mapsto-\mathbf r\).
A dipole operator \(\mathbf d=-e\mathbf r\)
changes sign. If the initial and final
states have identical parity, the
integrand of \(\langle f|\mathbf d|i\rangle\)
is odd and the integral vanishes.
Therefore a nonzero 3-j symbol does
\emph{not} by itself guarantee E1 emission.

\begin{checkpoint}
Use the shortest sequence:
(1) determine total \(J\) and projection
possibilities; (2) test 3-j triangle and
magnetic-number rules; (3) check parity;
(4) only then evaluate a 3-j coefficient
and reduced atomic matrix element.
\end{checkpoint}
